\begin{example}[5.6.11]
\label{IV.2.s5.6.11-fr}
Sea $A$ un anillo local noetheriano íntegro e \emph{integralmente cerrado}; entonces,
se ha visto (\textbf 0, 16.1.6) que la conclusión de (5.6.10) es válida para toda
$A$-álgebra finita íntegra $B$ que contiene a $A$. En cambio, vamos a
construir un ejemplo de anillo local noetheriano íntegro \emph{catenario} $A$ para el
cual no se cumplirá la conclusión de (5.6.10). Usaremos la construcción de
(5.2.5, (i)). Sean $k_0$ un cuerpo, $k$ una extensión puramente trascendente
$k_0(S_i)_{i\in\mathbb N}$ de grado de trascendencia infinito, $V$ el anillo
de valoración discreta $k[S]_{(S)}$, anillo local del anillo de polinomios
$k[S]$ en el ideal primo $(S)$; finalmente, sea $E$ el anillo de polinomios
$V[T]$; se ha visto que en $E$ el ideal maximal
$\mathfrak m=(S)+(T)$ es de altura 2 y el ideal maximal
$\mathfrak m'=(ST-1)$ de altura 1;

\oldpage[IV-2]{102}

los cuerpos residuales correspondientes son $\boldsymbol{k}(\mathfrak m)=k$ y
$\boldsymbol{k}(\mathfrak m')=k(S)$, cuerpo de fracciones de $V$; en virtud de la elección
de $k$, estos cuerpos son \emph{isomorfos}. Designemos por $\varepsilon$ y
$\varepsilon'$ los homomorfismos canónicos de $E$ sobre
$E/\mathfrak m=\boldsymbol{k}(\mathfrak m)$ y $E/\mathfrak m'=\boldsymbol{k}(\mathfrak m')$
respectivamente; sea, por otra parte, $\sigma$ un isomorfismo de
$\boldsymbol{k}(\mathfrak m)$ sobre $\boldsymbol{k}(\mathfrak m')$, y consideremos el subanillo $C$
de $E$ formado por los $x\in E$ tales que
$\varepsilon'(x)=\sigma(\varepsilon(x))$ (esta construcción es un caso
particular de los «procedimientos de pegado» que serán estudiados
sistemáticamente en el cap. V). Es inmediato que
$\mathfrak n=\mathfrak m\mathfrak m'=\mathfrak m\cap\mathfrak m'$ es un
ideal maximal de $C$, identificándose $C/\mathfrak m\mathfrak m'$ con el
subcuerpo de $(E/\mathfrak m)\times(E/\mathfrak m')$ formado por los pares
$(z,\sigma(z))$. Se tiene $E=C+C(ST-1)$; en otras palabras, $E$ es una
$C$-álgebra \emph{finita} y es evidentemente la clausura integral de $C$; veremos
además que $C$ es \emph{noetheriano}: esto se seguirá del lema siguiente:

